\subsection{Conjunto de datos y particiones}

Caltech Camera Traps (CCT-20), vía LILA BC: 243\,100 imágenes de 140 ubicaciones
de cámara, con 21 categorías de animales más la clase vacía. De esas, 63\,025
tienen bounding box, y solo esas pueden entrenar un detector: 46\,026 imágenes
únicas, 18\,GB, verificadas una por una contra truncamiento tras la descarga.

Se usa la partición de ECCV 2018~\cite{beery2018recognition}, que separa por
\emph{tipo de ubicación}: \texttt{train} (12\,099 imágenes con caja, 10
ubicaciones), \texttt{cis\_val} (1\,590) y \texttt{cis\_test} (12\,696) en las
\emph{mismas} ubicaciones, y \texttt{trans\_val} (1\,533) y \texttt{trans\_test}
(18\,033) en ubicaciones \emph{nunca vistas}.

\emph{Fuga por ráfagas.} Las cámaras disparan ráfagas de unas tres fotos casi
idénticas, y una ráfaga repartida entre entrenamiento y evaluación pone en la
evaluación una imagen que el modelo prácticamente memorizó. Medido:
\texttt{cis\_val} comparte 224 ráfagas con \texttt{train}, lo que contamina 75 de
sus 1\,665 imágenes con caja (4.5\%); esas 75 se excluyen. Las particiones
\emph{trans} y \texttt{cis\_test} no comparten ninguna, consistente con el
diseño: una ráfaga no cruza cámaras.

\emph{Las imágenes vacías no son una clase.} \texttt{empty} se excluye del mapa
de clases —y con ello del entrenamiento y de la evaluación, que usan solo
imágenes con caja—, y de las 21 categorías de animales del taxonómico de CCT solo
15 aparecen con bounding box en estas particiones, así que el mapa de clases —que
se construye desde las cajas presentes— queda en 15 especies. La razón de excluir
\texttt{empty} es conceptual: en detección
la ausencia de cajas ya expresa que no hay nada, y una clase de fondo le
enseñaría al modelo a dibujar una caja alrededor de la nada, con coordenadas que
no existen. Podrían sumarse como \emph{negativos explícitos} (una imagen sin
archivo de etiquetas cuenta como fondo); no se hizo porque habría cambiado a la
vez el tamaño del conjunto y la proporción fondo/objeto, y el efecto no sería
atribuible a ninguna de las dos cosas.

Dos conversiones de formato quedan fijadas por pruebas unitarias con un caso
verificado a mano, porque mal resueltas no producen ningún error visible: COCO
guarda la esquina superior izquierda más el tamaño en píxeles y el detector
espera el centro más el tamaño normalizado (confundirlos desplaza cada caja media
caja), y los identificadores de categoría de COCO no son contiguos mientras el
detector exige $0..N{-}1$ en el orden exacto de la lista de nombres.

\subsection{Decisiones de entrenamiento}
\label{sec:metodo_aumento}

\begin{table}[t]
  \centering
  \footnotesize
  \begin{tabular}{ll@{\qquad}ll}
    \toprule
    Decisión & Valor & Decisión & Valor \\
    \midrule
    Modelo base & \texttt{yolo11n.pt} & Volteo horizontal / vertical & 0.5 / 0.0 \\
    Resolución de entrada & 640 px & Mosaic (apagado en las últimas) & 1.0 (10 épocas) \\
    Épocas / lote & 60 / 16 & Escala / traslación & $\pm$50\% / $\pm$10\% \\
    Paciencia & 15 & HSV (h/s/v) & 0.015 / 0.5 / 0.4 \\
    Semilla & 42 & Rotación / perspectiva & 0.0 / 0.0 \\
    \bottomrule
  \end{tabular}
  \caption{Decisiones de entrenamiento (entregable \#1). El volteo vertical, la rotación y la perspectiva van en cero por la misma razón: la cámara está fija en un poste y esas vistas no existen en el despliegue. Escala y traslación son los valores por defecto de la implementación, declarados explícitamente porque están activos y pesan: la escala de $\pm$50\% es el aumento más relevante para el entregable \#4.}
  \label{tab:entrenamiento}
\end{table}


\emph{Por qué YOLO11\,n.} El enunciado sugiere YOLOv8 o YOLOv9; se usa YOLO11,
la iteración siguiente de la misma familia de detección en una sola pasada y la
versión vigente de Ultralytics, con la misma interfaz de entrenamiento y
evaluación. La variante \texttt{n}, la más pequeña, se eligió porque el factor
limitante del dominio es el número de ejemplos por especie y no la capacidad del
modelo, y porque permite un lote de 16 a 640 px en la GPU disponible; su costo en
los animales pequeños se discute entre los límites.

El criterio que gobierna el aumento de datos es uno solo: un aumento vale si
produce imágenes que la cámara podría haber tomado de verdad. La cámara está
atada a un poste y no gira, así que el volteo vertical, la rotación y la
perspectiva van en cero. El volteo horizontal sí vale, porque un animal cruza en
cualquier dirección y la especie no cambia al reflejarla, y la variación HSV
reproduce el rango real de iluminación, que incluye disparos nocturnos con
infrarrojo. \emph{Mosaic} compone cuatro imágenes en una: ayuda con objetos
chicos, que aquí son el caso dominante, y rompe el contexto de fondo, que es
justamente lo que el modelo tiende a memorizar. Se apaga en las últimas diez
épocas, porque terminar el entrenamiento con imágenes compuestas ajusta el modelo
a una distribución que no es la del despliegue.

\emph{Un aumento que el modelo nunca recibió.} El registro de la corrida
declara \texttt{erasing: 0.4} y \texttt{auto\_augment: randaugment}, y los dos
son engañosos: la implementación los aplica solo en clasificación. Verificarlo
importa para la validez del entregable \#4: el modelo cuya robustez ante oclusión
se mide en la Sección~\ref{sec:robustez} \emph{nunca vio oclusión durante el
entrenamiento}. Haber dado por buenas esas banderas habría invalidado la
interpretación del resultado.

Los hiperparámetros de optimización quedaron en los valores por defecto:
ajustarlos sin un conjunto de validación dedicado habría equivalido a ajustar a
\texttt{cis\_val}, que es una de las particiones que luego se reporta.

\subsection{Métricas, reimplementadas}
\label{sec:metricas}

Las métricas de detección se calculan con código propio por tres razones.

\emph{Clases sin positivos.} \texttt{deer} tiene 44 cajas de entrenamiento y
\emph{cero} en \texttt{trans\_test}. Su AP ahí no es cero: no está definida,
igual que un AUC sin positivos, y promediarla como cero castigaría al modelo con
una especie que era imposible de detectar porque no estaba. El código devuelve
indefinido y lista aparte las especies ausentes.

\emph{Subconjuntos arbitrarios.} Los entregables \#3 y \#4 piden evaluar lo
mismo por ubicación de cámara y por tamaño del objeto. Las predicciones se
guardan una vez (con umbral de confianza 0.001, el de COCO, porque recortar las
detecciones de baja confianza trunca la curva por el lado del recall alto e
\emph{infla} la métrica) y cada corte se recalcula sobre una lista filtrada en
segundos, en vez de una pasada completa del modelo por pregunta.

\emph{Incertidumbre.} El bootstrap remuestrea \emph{imágenes} y no cajas: las
cajas de una misma imagen comparten cámara, iluminación y a menudo el mismo
animal en ráfaga, y remuestrearlas por separado daría un intervalo demasiado
angosto.

La AP usa la interpolación de todos los puntos con envolvente monótona de PASCAL
VOC desde 2010~\cite{everingham2010pascal} y el emparejamiento codicioso por
confianza de COCO~\cite{lin2014microsoft}, sin reutilizar cajas anotadas. Se
reportan dos agregados: el \emph{macro}, que promedia las especies por igual y
por lo tanto expone el desbalance, y el \emph{ponderado por soporte} al lado;
la diferencia entre ambos \emph{es} la medida de ese desbalance. Como tercer
indicador, la correlación de rangos de Spearman entre el soporte de cada especie
y su AP; se usa Spearman y no Pearson porque el soporte va de 5 a 4\,800 cajas,
escala en la que Pearson queda dominada por las dos especies más comunes.
